\documentclass[11pt]{article}
\usepackage{amsmath,amssymb,amsthm,mathtools}
\usepackage{enumitem}
\usepackage{hyperref}
\usepackage[T1]{fontenc}
\title{Sharp Finite-Window Recognition and Information Loss in Oriented Pro-$p$ Groups}
\author{Seo Seongkyo}
\date{September 2026}
\newtheorem{theorem}{Theorem}[section]
\newtheorem{proposition}[theorem]{Proposition}
\newtheorem{lemma}[theorem]{Lemma}
\newtheorem{corollary}[theorem]{Corollary}
\newtheorem{remark}[theorem]{Remark}
\newtheorem{definition}[theorem]{Definition}
\begin{document}
\maketitle

\begin{abstract}
We study the exact Zassenhaus depth needed to factor affine crossed-cocycle representations at finite coefficient level. Let $A_k=\mathbb Z/p^k\mathbb Z$, $U_{1,k}=1+pA_k$, and $S_k=A_k\rtimes U_{1,k}$. For the standard Demushkin family
\[
G_{f,d}=\langle x_1,\ldots,x_d\mid x_1^{p^f}[x_1,x_2][x_3,x_4]\cdots[x_{d-1},x_d]=1\rangle,
\]
with odd $p$, $f\ge1$, and even $d\ge2$, every continuous affine representation factors through
$Q_k(G)=G/P_{p^{k-1}+1}(G)$, and this depth is sharp in the category of all affine crossed-cocycle representations:
$ n_{\mathrm{aff}}(k)=p^{k-1}+1$.
The lower bound includes rank two and does not require the orientation component to be surjective. We also separate factorization from recognition: a sharp factorization window does not by itself imply that a Kummerian predicate selects a unique orientation.
\end{abstract}

\section{Introduction}
The canonical Demushkin orientation is classical: Labute identifies it through the dualizing module, while the Kummerian/cyclotomic formulation is developed in later work. See \cite{Labute1967,EfratQuadrelli2019,QuadrelliWeigel2022}.
The present question is narrower. Given a finite coefficient level $k$, how much Zassenhaus information is required to preserve every affine crossed-cocycle representation into $S_k$? And does the resulting quotient necessarily contain enough information to recognize a unique orientation?

These are different questions. The first is a factorization problem in a specified representation category. The second is a recognition problem whose answer depends on the input category and predicate. The principal theorem of this paper is the exact and sharp factorization window; the recognition discussion is a logical boundary result, not a new positive recognition theorem.

For a fixed pro-$p$ group $G$, define
\[
n_{\mathrm{aff}}(k;G):=\min\{n\ge1:\text{every continuous }\psi:G\to S_k\text{ kills }P_n(G)\}.
\]
When $G=G_{f,d}$ we abbreviate this to $n_{\mathrm{aff}}(k)$.

\begin{theorem}[Sharp affine finite window]
Let $p$ be odd, $f\ge1$, $d\ge2$ even, and let $G_{f,d}$ be as above. For every $k\ge2$ and every continuous
\[
\psi=(z,\rho):G_{f,d}\longrightarrow S_k,
\qquad z\in Z^1(G_{f,d},A_k(\rho)),
\]
$\psi$ factors uniquely through
\[
Q_k(G_{f,d})=G_{f,d}/P_{p^{k-1}+1}(G_{f,d}).
\]
Moreover some such representation is nontrivial on $P_{p^{k-1}}(G_{f,d})$. Hence
\[
\boxed{n_{\mathrm{aff}}(k)=p^{k-1}+1.}
\]
\end{theorem}

\begin{remark}
This is category-relative minimality. No absolute minimality among arbitrary intrinsic carriers is claimed.
\end{remark}

\section{The affine target and the finite window}
Put
\[
A_k=\mathbb Z/p^k\mathbb Z, \qquad U_j=1+p^jA_k, \qquad S_k=A_k\rtimes U_1,
\]
with $(a,u)(b,v)=(a+ub,uv)$.

\begin{lemma}[Finite affine filtration]
For $n\ge1$, let $e(n)=\lceil\log_p n\rceil$, with $e(1)=0$. Then
\[
P_n(S_k)=p^{e(n)}A_k\rtimes U_{e(n)+1}.
\]
In particular
\[
P_{p^{k-1}}(S_k)=p^{k-1}A_k\ne1,
\qquad
P_{p^{k-1}+1}(S_k)=1.
\]
\end{lemma}

\begin{proof}
Use the Jennings--Lazard product for the Zassenhaus filtration
\[
P_n(S_k)=\prod_{ip^j\ge n}\gamma_i(S_k)^{p^j}.
\]
The commutator action gives
$\gamma_2(S_k)=pA_k$ and $\gamma_i(S_k)=p^{i-1}A_k$ for $i\ge2$.
Indeed $[1+p,a]=pa$, and induction gives the higher terms.
Principal-unit power estimates give $U_1^{p^j}=U_{j+1}$, while
$(a,1)^{p^j}=(p^ja,1)$, hence
$S_k^{p^j}=p^jA_k\rtimes U_{j+1}$.
If $e=e(n)$, the $i=1,j=e$ term already supplies $p^eA_k\rtimes U_{e+1}$.
For $i\ge2$ with $ip^j\ge n$, one has $i-1+j\ge e$, hence
$\gamma_i(S_k)^{p^j}\subseteq p^eA_k$. This proves the formula and the endpoint values.
\end{proof}

\begin{proposition}[Arbitrary-candidate factorization]
For every pro-$p$ group $G$, every continuous $\rho:G\to U_{1,k}$, and every continuous $z\in Z^1(G,A_k(\rho))$, the homomorphism
\[
\psi(g)=(z(g),\rho(g))
\]
kills $P_{p^{k-1}+1}(G)$. Thus the crossed cocycle, its orientation component, and the associated twisted $H^1$ lifting problem factor through $Q_k(G)$.
\end{proposition}

\begin{proof}
The crossed-cocycle identity is exactly the homomorphism condition into $S_k$. Zassenhaus filtrations are functorial, so
$\psi(P_{p^{k-1}+1}(G))\subseteq P_{p^{k-1}+1}(S_k)=1$.
Thus both the coefficient action and the cocycle factor through the quotient, and inflation identifies the corresponding cocycles and coboundaries with their quotient-level counterparts.
\end{proof}

\section{Sharpness}
\begin{proposition}[The case $f<k$]
Let $f<k$. For the canonical orientation,
\[
\chi(x_1)=\chi(x_3)=\cdots=\chi(x_d)=1,
\qquad
\chi(x_2)=(1-p^f)^{-1}\pmod{p^k}.
\]
Set $z(x_1)=1$ and $z(x_i)=0$ for $i>1$. Then $z$ descends to a crossed cocycle and
\[
z(x_1^{p^{k-1}})=p^{k-1}\ne0\pmod{p^k}.
\]
\end{proposition}

\begin{proof}
With the convention $z(gh)=z(g)+\chi(g)z(h)$, the power--commutator part of the relator evaluates to
\[
p^f+\chi(x_2)^{-1}(1-\chi(x_2))=0,
\]
because $\chi(x_2)=(1-p^f)^{-1}$. The remaining commutators have zero cocycle value. Since $\chi(x_1)=1$, $z(x_1^m)=m$.
\end{proof}

\begin{proposition}[The case $f\ge k$, including rank two]
Let $f\ge k$ and define
\[
\rho(x_1)=1,\quad \rho(x_2)=1+p,\quad z(x_1)=0,\quad z(x_2)=1,
\]
with all unused generator values and cocycle values trivial. Then $z$ is a crossed cocycle and
\[
z(x_2^{p^{k-1}})
= \frac{(1+p)^{p^{k-1}}-1}{p}\ne0\pmod{p^k}.
\]
\end{proposition}

\begin{proof}
The term $x_1^{p^f}$ is invisible at level $p^k$, and $z(x_1)=0$ kills the commutator contribution. The cyclic cocycle formula gives the displayed geometric sum. LTE yields
$v_p((1+p)^{p^{k-1}}-1)=k$, so the quotient has valuation $k-1$.
Only $x_1,x_2$ are used, so the construction covers $d=2$.
\end{proof}

\begin{theorem}[Category-relative minimality]
For every odd $p$, every $f\ge1$, every even $d\ge2$, and every $k\ge2$,
\[
\boxed{n_{\mathrm{aff}}(k)=p^{k-1}+1}
\]
for the category of all continuous crossed-cocycle representations into $A_k(\rho)\rtimes U_{1,k}$.
\end{theorem}

\begin{proof}
The filtration lemma gives the upper bound. The two preceding propositions give a witness nontrivial on $P_{p^{k-1}}$ in the cases $f<k$ and $f\ge k$, respectively. Thus no smaller depth works for the entire affine category.
\end{proof}

\begin{remark}
No surjectivity assumption on $\rho(G)$ is needed. For $f<k$ the canonical orientation may fail to generate $U_{1,k}$, but the translation component already detects the missing Zassenhaus layer.
\end{remark}

\section{Finite-depth q-collapse}
Let
\[
Q_k^{(f)}=G_{f,d}/P_{p^{k-1}+1}(G_{f,d}).
\]
If $f\ge k$, then $x_1^{p^f}\in P_{p^f}(G_{f,d})$, and $p^f\ge p^k>p^{k-1}+1$. Hence the finite quotient has the same defining relation as the $f=\infty$ relation
\[
[x_1,x_2][x_3,x_4]\cdots[x_{d-1},x_d]=1.
\]
Thus all $f\ge k$ collapse to the same finite quotient after identifying the generators. At the same time
$(1-p^f)^{-1}\equiv1\pmod{p^k}$, so the canonical orientation has the same reduction.
For $f<k$, one has $G_{f,d}^{\mathrm{ab}}\simeq \mathbb Z_p^{d-1}\oplus\mathbb Z/p^f\mathbb Z$, while quotienting by $P_{p^{k-1}+1}$ reduces the free abelian directions only modulo $p^k$. Hence the image of $[x_1]$ still has exact order $p^f$ in the abelianization of $Q_k^{(f)}$, so the finite torsion order remains visible.

\section{Factorization is not recognition}
For a continuous candidate $\rho:Q_k\to U_{1,k}$ define
\[
\mathsf K_k(Q_k,\rho):\quad
H^1(Q_k,A_k(\rho))\longrightarrow H^1(Q_k,\mathbb F_p)
\quad\text{is surjective}.
\]
The predicate is natural in $(Q_k,\rho)$ and has no $f$ or presentation as input. The factorization theorem says that the corresponding twisted lifting problem factors through $Q_k$ for every candidate.

However, uniqueness is an additional property of the input class.

\begin{proposition}[Information loss in free pro-$p$ groups]
Let $F_d$ be free pro-$p$ of finite rank $d\ge2$. For every continuous orientation $\rho:F_d\to U_{1,k}$, the finite-level lifting map
\[
H^1(F_d,A_k(\rho))\longrightarrow H^1(F_d,\mathbb F_p)
\]
is surjective. Thus the Kummerian predicate itself does not single out a unique orientation on the class of free pro-$p$ groups.
\end{proposition}

\begin{proof}
$\operatorname{cd}_p(F_d)=1$, so $H^2(F_d,M)=0$ for every finite discrete $p$-primary module $M$. Therefore the connecting map for
$0\to pA_k(\rho)\to A_k(\rho)\to\mathbb F_p\to0$
vanishes and the $H^1$ reduction map is surjective. Since $F_d$ has many distinct orientations, the predicate itself cannot be a uniqueness selector on this class.
\end{proof}

\begin{remark}
This does not prove that every conceivable isomorphism-natural selector is impossible. Such a claim would require a separate automorphism or indistinguishability argument for the exact quotient category.
\end{remark}

\section{Finite free products of Demushkin blocks and applications}
Let
\[
G=D_1*_p\cdots *_pD_m
\]
be a finite free pro-$p$ product of Demushkin groups. This section records concrete consequences of the sharp affine-window theorem. The upper-bound statement is intrinsic and does not require chosen presentations or orientations on the factors; the lower-bound and parameter-profile statements use the standard Demushkin presentations only where explicitly indicated.

Write
\[
T_n(G)=G/P_n(G).
\]

\begin{lemma}[Zassenhaus truncation and coproducts]
For pro-$p$ groups $G_1,G_2$,
\[
T_n(G_1*_pG_2)
\cong
\frac{T_n(G_1)*_pT_n(G_2)}
{P_n(T_n(G_1)*_pT_n(G_2))}.
\]
\end{lemma}

\begin{proof}
$T_n$ is left adjoint to the inclusion of pro-$p$ groups $H$ satisfying
$P_n(H)=1$, because every continuous map $G\to H$ kills $P_n(G)$.
Left adjoints preserve coproducts, giving the formula.
\end{proof}

\begin{theorem}[Uniform affine window for finite Demushkin free products]
For a finite free pro-$p$ product of Demushkin groups, every continuous affine crossed-cocycle representation at level $k$ factors through
\[
G/P_{p^{k-1}+1}(G).
\]
Thus the same depth is sufficient uniformly for this class. If one factor is a standard $G_{f,d}$, the sharp witness in the preceding section shows that no smaller depth works uniformly for the whole class.
\end{theorem}

\begin{proof}
The factorization follows directly from the arbitrary-candidate factorization proposition. The truncation lemma shows that no separate mixed-commutator window is needed: mixed words are absorbed into the reflected $P_{p^{k-1}+1}$-kernel. For the lower bound, take a sharp affine witness on one factor and compose it with the canonical factor inclusion into the free pro-$p$ coproduct.
\end{proof}

\subsection{Finite-depth parameter collapse}

For $f\ge1$ and even $d\ge2$, consider the standard block
\[
D_{f,d}=
\left\langle x_1,\ldots,x_d\ \middle|\
x_1^{p^f}[x_1,x_2]\cdots[x_{d-1},x_d]=1
\right\rangle .
\]
The power term has initial Zassenhaus weight $p^f$:
\[
x_1^{p^f}\in P_{p^f}\setminus P_{p^f+1}
\]
in the free pro-$p$ group before imposing the relation. Hence at the window
\[
Q_k(D_{f,d})=D_{f,d}/P_{p^{k-1}+1}(D_{f,d})
\]
the power term is invisible whenever $f\ge k$.

\begin{proposition}[Finite-depth $f$-collapse]
For the standard block $D_{f,d}$, if $f\ge k$, the image of the defining relation in $Q_k(D_{f,d})$ agrees with the power-free relation
\[
[x_1,x_2]\cdots[x_{d-1},x_d]=1.
\]
If $f<k$, the torsion class represented by $x_1$ survives with order $p^f$ in the abelianization of the finite quotient. Thus the marked finite quotient distinguishes the two regimes
\[
f<k,\qquad f\ge k.
\]
\end{proposition}

\begin{proof}
If $f\ge k$, then $p^f\ge p^k>p^{k-1}+1$, so
$x_1^{p^f}\in P_{p^{k-1}+1}$ and disappears in the quotient.
For $f<k$, the abelianization of the standard block is
\[
D_{f,d}^{\mathrm{ab}}\cong\mathbb Z_p^{\,d-1}\oplus\mathbb Z/p^f\mathbb Z.
\]
The image of $P_{p^{k-1}+1}$ cannot remove the pre-existing order-$p^f$
torsion class, since $f<k$ and the finite window only imposes relations
of $p$-power order at least $p^k$ in the corresponding free directions.
Hence the marked abelianization retains the order-$p^f$ class.
\end{proof}

\begin{corollary}[Heterogeneous parameter profile]
For
\[
G=*_{i=1}^mD_{f_i,d_i},
\]
the finite quotient $Q_k(G)$ processes all blocks at the same Zassenhaus
depth. In the inherited block marking, it separates the parameters with
$f_i<k$ from those with $f_i\ge k$: the latter have already entered the
power-free finite-depth regime, while the former retain their
$p^{f_i}$ torsion class in abelianization.
\end{corollary}

This is an application-level information-loss statement. It does not claim
that the bare abstract quotient canonically recovers the individual block
decomposition or that the quotient is an absolute minimal carrier.

\subsection{A heterogeneous composite sharpness example}

Consider
\[
G=D_{1,4}*_pD_{3,2}*_pD_{7,4}.
\]
At coefficient level $k=4$ the three parameters satisfy
\[
1<4,\qquad 3<4,\qquad 7\ge4.
\]
Thus the first two blocks retain their finite-depth power parameters,
whereas the third has already entered the power-free regime. Nevertheless
the entire composite group uses the single affine window
\[
Q_4(G)=G/P_{p^3+1}(G).
\]

The sharpness theorem remains visible on this composite object. Choose a
sharp witness on the first block and extend it trivially across the other
free factors. Then the resulting affine representation is still nontrivial
on $P_{p^3}(G)$, so the depth $p^3+1$ cannot be lowered uniformly for the
class containing this example.

The point is not a new sharpness theorem: the exact threshold is proved
above for the affine category. The example shows concretely that the same
sharp window survives a genuinely composite object with heterogeneous
Demu\v{s}kin parameters and mixed commutators.

\subsection{Blockwise finite Kummer recognition}

We now record the application to the recognition theorem of the preceding
paper, and restrict it exactly to the class for which that theorem is
available. Let
\[
G=G_1*_3\cdots *_3G_m
\]
where every $G_i$ is a rank-four $q=3$ Demu\v{s}kin block. For each factor,
let $\chi_i$ denote its canonical orientation. Paper 1 establishes, for
arbitrary principal-unit candidates $\rho_i$,
\[
\mathsf K_k(Q_k(G_i),\rho_i)
\Longleftrightarrow
\rho_i=\chi_i\pmod{3^k}.
\]

For a candidate
\[
\rho:Q_k(G)\longrightarrow U_{1,k},
\]
write $\rho_i$ for its restriction to the finite quotient associated to the
$i$th factor.

\begin{proposition}[Degree-one free-product decomposition]
For a finite coefficient module $M$ and a free pro-$p$ product,
\[
H^1(G,M)\cong\bigoplus_{i=1}^m H^1(G_i,M|_{G_i}),
\]
and the same decomposition holds with $M=\mathbb F_p$.
\end{proposition}

\begin{proof}
A continuous crossed homomorphism on a free pro-$p$ product is uniquely
determined by its restrictions to the free factors. Principal crossed
homomorphisms decompose compatibly, giving the stated degree-one
Mayer--Vietoris decomposition.
\end{proof}

\begin{theorem}[Blockwise finite Kummer recognition]
Under the preceding hypotheses,
\[
\mathsf K_k(Q_k(G),\rho)
\Longleftrightarrow
\mathsf K_k(Q_k(G_i),\rho_i)
\quad\text{for every }i.
\]
Consequently,
\[
\mathsf K_k(Q_k(G),\rho)
\Longleftrightarrow
\rho_i=\chi_i\pmod{3^k}
\quad\text{for every }i.
\]
\end{theorem}

\begin{proof}
By the truncation lemma, the finite quotient is the reflected
$P_{p^{k-1}+1}$-quotient of the free pro-$3$ product of the factor
truncations. The arbitrary-candidate factorization theorem identifies the
finite-level twisted $H^1$ problem on $G$ with the corresponding problem on
$Q_k(G)$ for every candidate. For trivial coefficients,
$P_{p^{k-1}+1}(G)\subseteq\Phi(G)$ gives the analogous $H^1$ identification.
The degree-one free-product decomposition therefore identifies the global
reduction map with the direct sum of the factor reduction maps. Such a
direct sum is surjective exactly when every factor map is surjective.
Applying the rank-four $q=3$ selector theorem from Paper 1 to each factor
gives the final equivalence.
\end{proof}

\begin{remark}
This is an application of the preceding recognition theorem, not a new
classification of orientations of arbitrary free products. In particular,
the theorem does not assert that a general free product is itself
Demu\v{s}kin or has a single canonical orientation. It recognizes the
canonical orientations of the declared factors block by block.
\end{remark}

\begin{remark}
No claim is made for broader recursively defined elementary-type classes
involving fibre products or amalgamated operations. Likewise, no absolute
minimality of $Q_k$ among arbitrary intrinsic carriers is asserted.
\end{remark}

\section{Relation with the preceding paper}
The preceding paper treats finite-window recognition for the fixed rank-four $p=3$ Demushkin group. Its intrinsic Kummer predicate on the corresponding finite quotient selects the canonical orientation modulo $3^k$. The present paper adds the exact factorization depth for the full affine crossed-cocycle category, extends sharpness to the rank-two boundary, records the finite-depth $f$-collapse, and separates factorization from recognition.

The logical chain is
\[
\text{affine factorization}
\quad\ne\quad
\quad\text{orientation recognition}
\quad\ne\quad
\quad\text{absolute information sufficiency}.
\]

\section{Literature boundary and novelty}
The canonical Demushkin orientation and its Kummerian characterization are classical \cite{Labute1967,EfratQuadrelli2019}. The general functoriality of the Zassenhaus filtration under continuous homomorphisms, and its representation-theoretic use, are standard; see Lazard/Jennings theory and Efrat's treatment of Zassenhaus filtrations and representations \cite{Efrat2014,MinacRogelstadTan2016}. Oriented elementary-type constructions and related Kummerian closure phenomena are treated by Quadrelli--Weigel \cite{QuadrelliWeigel2022}. Quadrelli's 2024 work studies 1-cyclotomicity and quotient-related obstructions \cite{Quadrelli2024}. These are surrounding results, not claims of prior identity with the present theorem. In particular, the literature establishes the standard Kummerian/1-cyclotomic framework and substantial information on Zassenhaus quotients, but the audit did not locate an earlier theorem stated in the present affine target category with the exact depth $p^{k-1}+1$ and the present sharpness witnesses.

The present contribution is deliberately narrower: the exact depth $p^{k-1}+1$ for factorization of all affine crossed-cocycle representations into $A_k\rtimes U_{1,k}$, explicit lower-bound witnesses through rank two, finite-depth $f$-collapse, and the separation of sharp factorization from uniqueness of a Kummerian orientation.

The novelty claim is conditional. The audited corpus did not reveal a theorem stated in exactly this representation category with exactly this depth and lower-bound formulation. This is a publication boundary statement, not an absolute priority claim.

\section{Limitations and conclusion}
The sharpness theorem is category-relative. The positive free-product theorem is restricted to finite free products of Demushkin blocks. The free-pro-$p$ example concerns the Kummer predicate itself, not every conceivable selector. The finite-depth q-collapse does not imply that the full groups with different $f$ are isomorphic.

The main formula is
\[
\boxed{n_{\mathrm{aff}}(k)=p^{k-1}+1.}
\]
It is uniform for the standard Demushkin parameters $f\ge1$ and even rank $d\ge2$, including rank two. Any stronger claim of absolute minimality or universal orientation reconstruction requires a separately specified carrier category and recognition theorem.

\begin{thebibliography}{99}
\bibitem{Efrat2014}
I. Efrat, \emph{The Zassenhaus filtration, Massey products, and representations of profinite groups}, Adv. Math. 263 (2014), 389--411.
\bibitem{MinacRogelstadTan2016}
J. Mináč, M. Rogelstad and N. D. Tân, \emph{Dimensions of Zassenhaus filtration subquotients of some pro-$p$-groups}, Israel J. Math. 212 (2016), 825--855.
\bibitem{Labute1967}
J. Labute, \emph{Classification of Demushkin Groups}, Canadian J. Math. 19 (1967), 106--132.
\bibitem{EfratQuadrelli2019}
I. Efrat and C. Quadrelli, \emph{The Kummerian Property and Maximal Pro-$p$ Galois Groups}, J. Algebra 525 (2019), 284--310.
\bibitem{QuadrelliWeigel2022}
C. Quadrelli and T. S. Weigel, \emph{Oriented pro-$\ell$ groups with the Bogomolov--Positselski property}, Res. Number Theory 8 (2022), 21.
\bibitem{Quadrelli2024}
C. Quadrelli, \emph{Chasing Maximal Pro-$p$ Galois Groups via 1-Cyclotomicity}, Mediterr. J. Math. 21 (2024), 56.
\end{thebibliography}
\end{document}